% Options for packages loaded elsewhere
% Options for packages loaded elsewhere
\PassOptionsToPackage{unicode}{hyperref}
\PassOptionsToPackage{hyphens}{url}
%
\documentclass[
  english,
  russian,
  12pt,
  a4paper,
  DIV=11,
  numbers=noendperiod]{scrreprt}
\usepackage{xcolor}
\usepackage{amsmath,amssymb}
\setcounter{secnumdepth}{5}
\usepackage{iftex}
\ifPDFTeX
  \usepackage[T1]{fontenc}
  \usepackage[utf8]{inputenc}
  \usepackage{textcomp} % provide euro and other symbols
\else % if luatex or xetex
  \usepackage{unicode-math} % this also loads fontspec
  \defaultfontfeatures{Scale=MatchLowercase}
  \defaultfontfeatures[\rmfamily]{Ligatures=TeX,Scale=1}
\fi
\usepackage{lmodern}
\ifPDFTeX\else
  % xetex/luatex font selection
\fi
% Use upquote if available, for straight quotes in verbatim environments
\IfFileExists{upquote.sty}{\usepackage{upquote}}{}
\IfFileExists{microtype.sty}{% use microtype if available
  \usepackage[]{microtype}
  \UseMicrotypeSet[protrusion]{basicmath} % disable protrusion for tt fonts
}{}
\usepackage{setspace}
% Make \paragraph and \subparagraph free-standing
\makeatletter
\ifx\paragraph\undefined\else
  \let\oldparagraph\paragraph
  \renewcommand{\paragraph}{
    \@ifstar
      \xxxParagraphStar
      \xxxParagraphNoStar
  }
  \newcommand{\xxxParagraphStar}[1]{\oldparagraph*{#1}\mbox{}}
  \newcommand{\xxxParagraphNoStar}[1]{\oldparagraph{#1}\mbox{}}
\fi
\ifx\subparagraph\undefined\else
  \let\oldsubparagraph\subparagraph
  \renewcommand{\subparagraph}{
    \@ifstar
      \xxxSubParagraphStar
      \xxxSubParagraphNoStar
  }
  \newcommand{\xxxSubParagraphStar}[1]{\oldsubparagraph*{#1}\mbox{}}
  \newcommand{\xxxSubParagraphNoStar}[1]{\oldsubparagraph{#1}\mbox{}}
\fi
\makeatother


\usepackage{longtable,booktabs,array}
\usepackage{calc} % for calculating minipage widths
% Correct order of tables after \paragraph or \subparagraph
\usepackage{etoolbox}
\makeatletter
\patchcmd\longtable{\par}{\if@noskipsec\mbox{}\fi\par}{}{}
\makeatother
% Allow footnotes in longtable head/foot
\IfFileExists{footnotehyper.sty}{\usepackage{footnotehyper}}{\usepackage{footnote}}
\makesavenoteenv{longtable}
\usepackage{graphicx}
\makeatletter
\newsavebox\pandoc@box
\newcommand*\pandocbounded[1]{% scales image to fit in text height/width
  \sbox\pandoc@box{#1}%
  \Gscale@div\@tempa{\textheight}{\dimexpr\ht\pandoc@box+\dp\pandoc@box\relax}%
  \Gscale@div\@tempb{\linewidth}{\wd\pandoc@box}%
  \ifdim\@tempb\p@<\@tempa\p@\let\@tempa\@tempb\fi% select the smaller of both
  \ifdim\@tempa\p@<\p@\scalebox{\@tempa}{\usebox\pandoc@box}%
  \else\usebox{\pandoc@box}%
  \fi%
}
% Set default figure placement to htbp
\def\fps@figure{htbp}
\makeatother



\ifLuaTeX
\usepackage[bidi=basic,provide=*]{babel}
\else
\usepackage[bidi=default,provide=*]{babel}
\fi
% get rid of language-specific shorthands (see #6817):
\let\LanguageShortHands\languageshorthands
\def\languageshorthands#1{}


\setlength{\emergencystretch}{3em} % prevent overfull lines

\providecommand{\tightlist}{%
  \setlength{\itemsep}{0pt}\setlength{\parskip}{0pt}}



 
\usepackage[backend=biber,langhook=extras,autolang=other*]{biblatex}
\addbibresource{bib/cite.bib}

\usepackage[]{csquotes}

\KOMAoption{captions}{tableheading}
\usepackage{indentfirst}
\usepackage{float}
\floatplacement{figure}{H}
\usepackage{libertine}
\makeatletter
\@ifpackageloaded{caption}{}{\usepackage{caption}}
\AtBeginDocument{%
\ifdefined\contentsname
  \renewcommand*\contentsname{Содержание}
\else
  \newcommand\contentsname{Содержание}
\fi
\ifdefined\listfigurename
  \renewcommand*\listfigurename{Список иллюстраций}
\else
  \newcommand\listfigurename{Список иллюстраций}
\fi
\ifdefined\listtablename
  \renewcommand*\listtablename{Список таблиц}
\else
  \newcommand\listtablename{Список таблиц}
\fi
\ifdefined\figurename
  \renewcommand*\figurename{Рисунок}
\else
  \newcommand\figurename{Рисунок}
\fi
\ifdefined\tablename
  \renewcommand*\tablename{Таблица}
\else
  \newcommand\tablename{Таблица}
\fi
}
\@ifpackageloaded{float}{}{\usepackage{float}}
\floatstyle{ruled}
\@ifundefined{c@chapter}{\newfloat{codelisting}{h}{lop}}{\newfloat{codelisting}{h}{lop}[chapter]}
\floatname{codelisting}{Список}
\newcommand*\listoflistings{\listof{codelisting}{Листинги}}
\makeatother
\makeatletter
\makeatother
\makeatletter
\@ifpackageloaded{caption}{}{\usepackage{caption}}
\@ifpackageloaded{subcaption}{}{\usepackage{subcaption}}
\makeatother
\usepackage{bookmark}
\IfFileExists{xurl.sty}{\usepackage{xurl}}{} % add URL line breaks if available
\urlstyle{same}
\hypersetup{
  pdftitle={Отчёт для первого модуля внешнего курса},
  pdfauthor={Семенченко Татьяна Сергеевна},
  pdflang={ru-RU},
  hidelinks,
  pdfcreator={LaTeX via pandoc}}


\title{Отчёт для первого модуля внешнего курса}
\usepackage{etoolbox}
\makeatletter
\providecommand{\subtitle}[1]{% add subtitle to \maketitle
  \apptocmd{\@title}{\par {\large #1 \par}}{}{}
}
\makeatother
\subtitle{Модуль 1}
\author{Семенченко Татьяна Сергеевна}
\date{}
\begin{document}
\maketitle

\renewcommand*\contentsname{Содержание}
{
\setcounter{tocdepth}{1}
\tableofcontents
}
\listoffigures
\listoftables

\setstretch{1.5}
\chapter{Цель
работы}\label{ux446ux435ux43bux44c-ux440ux430ux431ux43eux442ux44b}

Пройти внешний курс \enquote{Введение в Linux!}

\chapter{Задание}\label{ux437ux430ux434ux430ux43dux438ux435}

Прохождение 1 модуля внешнего курса на stepic \enquote{Введение в
Linux!}.

\chapter{Выполнение
курса}\label{ux432ux44bux43fux43eux43bux43dux435ux43dux438ux435-ux43aux443ux440ux441ux430}

\section{Общая информация о
курсе}\label{ux43eux431ux449ux430ux44f-ux438ux43dux444ux43eux440ux43cux430ux446ux438ux44f-ux43e-ux43aux443ux440ux441ux435}

Курс называется \enquote{Введение в Linux!}(рис.~\ref{fig-01}).

\begin{figure}

\centering{

\pandocbounded{\includegraphics[keepaspectratio]{image/1.png}}

}

\caption{\label{fig-01}}

\end{figure}%

Правильные ответы в вопросе прохождения курса (рис.~\ref{fig-02}).

\begin{figure}

\centering{

\pandocbounded{\includegraphics[keepaspectratio]{image/2.png}}

}

\caption{\label{fig-02}}

\end{figure}%

\section{Как установить
Linux}\label{ux43aux430ux43a-ux443ux441ux442ux430ux43dux43eux432ux438ux442ux44c-linux}

Я использую windows операционную систему (рис.~\ref{fig-03}).

\begin{figure}

\centering{

\pandocbounded{\includegraphics[keepaspectratio]{image/3.png}}

}

\caption{\label{fig-03}}

\end{figure}%

Виртуальная машина - это специальная программа для запуска одной ОС на
другой ОС (рис.~\ref{fig-04})

\begin{figure}

\centering{

\pandocbounded{\includegraphics[keepaspectratio]{image/4.png}}

}

\caption{\label{fig-04}}

\end{figure}%

Да, я работаю с системе Linux (рис.~\ref{fig-05}).

\begin{figure}

\centering{

\pandocbounded{\includegraphics[keepaspectratio]{image/5.png}}

}

\caption{\label{fig-05}}

\end{figure}%

\section{Осваиваем
Linux}\label{ux43eux441ux432ux430ux438ux432ux430ux435ux43c-linux}

Создала документ LibreOffice Writer и написала шрифтом FreeMono Hello,
Linux! (рис.~\ref{fig-32}).

\begin{figure}

\centering{

\pandocbounded{\includegraphics[keepaspectratio]{image/32.png}}

}

\caption{\label{fig-32}}

\end{figure}%

Ubuntu имеет расширение deb (рис.~\ref{fig-06}).

\begin{figure}

\centering{

\pandocbounded{\includegraphics[keepaspectratio]{image/6.png}}

}

\caption{\label{fig-06}}

\end{figure}%

Установила VLC и написала фамилию \enquote{Denis-Courmont}
(рис.~\ref{fig-07}).

\begin{figure}

\centering{

\pandocbounded{\includegraphics[keepaspectratio]{image/7.png}}

}

\caption{\label{fig-07}}

\end{figure}%

Update Manager можно использовать для обновления ссылок в Software
Center, установленных программ, всей системы до новой версии
(рис.~\ref{fig-08}).

\begin{figure}

\centering{

\pandocbounded{\includegraphics[keepaspectratio]{image/8.png}}

}

\caption{\label{fig-08}}

\end{figure}%

\section{Terminal:
основы}\label{terminal-ux43eux441ux43dux43eux432ux44b}

Синонимами являются Терминал и Консоль (рис.~\ref{fig-09}).

\begin{figure}

\centering{

\pandocbounded{\includegraphics[keepaspectratio]{image/9.png}}

}

\caption{\label{fig-09}}

\end{figure}%

Команда pwd (рис.~\ref{fig-10}).

\begin{figure}

\centering{

\pandocbounded{\includegraphics[keepaspectratio]{image/10.png}}

}

\caption{\label{fig-10}}

\end{figure}%

Указала верные команды (рис.~\ref{fig-11}).

\begin{figure}

\centering{

\pandocbounded{\includegraphics[keepaspectratio]{image/11.png}}

}

\caption{\label{fig-11}}

\end{figure}%

Указала верные команды (рис.~\ref{fig-12}).

\begin{figure}

\centering{

\pandocbounded{\includegraphics[keepaspectratio]{image/12.png}}

}

\caption{\label{fig-12}}

\end{figure}%

Команда rm -r (рис.~\ref{fig-13}).

\begin{figure}

\centering{

\pandocbounded{\includegraphics[keepaspectratio]{image/13.png}}

}

\caption{\label{fig-13}}

\end{figure}%

\section{Запуск исполняемых
файлов}\label{ux437ux430ux43fux443ux441ux43a-ux438ux441ux43fux43eux43bux43dux44fux435ux43cux44bux445-ux444ux430ux439ux43bux43eux432}

После команды Exit ничего не закроется (рис.~\ref{fig-14}).

\begin{figure}

\centering{

\pandocbounded{\includegraphics[keepaspectratio]{image/14.png}}

}

\caption{\label{fig-14}}

\end{figure}%

Запуск, Ctrl+Z, bg (рис.~\ref{fig-15}).

\begin{figure}

\centering{

\pandocbounded{\includegraphics[keepaspectratio]{image/15.png}}

}

\caption{\label{fig-15}}

\end{figure}%

Исполнила файл через python3, скопировала его вывод (рис.~\ref{fig-16}).

\begin{figure}

\centering{

\pandocbounded{\includegraphics[keepaspectratio]{image/16.png}}

}

\caption{\label{fig-16}}

\end{figure}%

\begin{figure}

\centering{

\pandocbounded{\includegraphics[keepaspectratio]{image/17.png}}

}

\caption{\label{fig-17}}

\end{figure}%

\section{Ввод /
вывод}\label{ux432ux432ux43eux434-ux432ux44bux432ux43eux434}

Поток ошибок выводится на экран (рис.~\ref{fig-18}).

\begin{figure}

\centering{

\pandocbounded{\includegraphics[keepaspectratio]{image/18.png}}

}

\caption{\label{fig-18}}

\end{figure}%

Выбрала команды с \textgreater, \textgreater\textgreater{}
(рис.~\ref{fig-20}).

\begin{figure}

\centering{

\pandocbounded{\includegraphics[keepaspectratio]{image/20.png}}

}

\caption{\label{fig-20}}

\end{figure}%

Сообщения об ошибках выводятся на экран (рис.~\ref{fig-21}).

\begin{figure}

\centering{

\pandocbounded{\includegraphics[keepaspectratio]{image/21.png}}

}

\caption{\label{fig-21}}

\end{figure}%

\section{Скачивание файлов из
интернета}\label{ux441ux43aux430ux447ux438ux432ux430ux43dux438ux435-ux444ux430ux439ux43bux43eux432-ux438ux437-ux438ux43dux442ux435ux440ux43dux435ux442ux430}

Картинка окажется в /home/alex/1.jpg (рис.~\ref{fig-22}).

\begin{figure}

\centering{

\pandocbounded{\includegraphics[keepaspectratio]{image/22.png}}

}

\caption{\label{fig-22}}

\end{figure}%

Опцию -q или --quiet (рис.~\ref{fig-23}).

\begin{figure}

\centering{

\pandocbounded{\includegraphics[keepaspectratio]{image/23.png}}

}

\caption{\label{fig-23}}

\end{figure}%

Выбрала ответ: Будут скачаны jpg и html файлы, но все html будут удалены
(рис.~\ref{fig-24}).

\begin{figure}

\centering{

\pandocbounded{\includegraphics[keepaspectratio]{image/24.png}}

}

\caption{\label{fig-24}}

\end{figure}%

\section{Работа с
архивами}\label{ux440ux430ux431ux43eux442ux430-ux441-ux430ux440ux445ux438ux432ux430ux43cux438}

gzip удаляет архив после его распаковки (рис.~\ref{fig-25}).

\begin{figure}

\centering{

\pandocbounded{\includegraphics[keepaspectratio]{image/25.png}}

}

\caption{\label{fig-25}}

\end{figure}%

zip и tar (рис.~\ref{fig-26}).

\begin{figure}

\centering{

\pandocbounded{\includegraphics[keepaspectratio]{image/26.png}}

}

\caption{\label{fig-26}}

\end{figure}%

Нужно указать -cjf (рис.~\ref{fig-27}).

\begin{figure}

\centering{

\pandocbounded{\includegraphics[keepaspectratio]{image/27.png}}

}

\caption{\label{fig-27}}

\end{figure}%

\section{Поиск файлов и слов в
файлах}\label{ux43fux43eux438ux441ux43a-ux444ux430ux439ux43bux43eux432-ux438-ux441ux43bux43eux432-ux432-ux444ux430ux439ux43bux430ux445}

Команды \emph{.jpg, alexey.}, *.? НЕ найдут файл Alexey.jpeg
(рис.~\ref{fig-28}).

\begin{figure}

\centering{

\pandocbounded{\includegraphics[keepaspectratio]{image/28.png}}

}

\caption{\label{fig-28}}

\end{figure}%

Выбираем там, где есть world учитывая регистр, должно быть с маленькой
буквы (рис.~\ref{fig-29}).

\begin{figure}

\centering{

\pandocbounded{\includegraphics[keepaspectratio]{image/29.png}}

}

\caption{\label{fig-29}}

\end{figure}%

Скачала архив, сгенерировала и загрузила файл в форму
(рис.~\ref{fig-30}).

\begin{figure}

\centering{

\pandocbounded{\includegraphics[keepaspectratio]{image/30.png}}

}

\caption{\label{fig-30}}

\end{figure}%

\begin{figure}

\centering{

\pandocbounded{\includegraphics[keepaspectratio]{image/31.png}}

}

\caption{\label{fig-31}}

\end{figure}%

\chapter{Выводы}\label{ux432ux44bux432ux43eux434ux44b}

Прошёла Модуль 1 курса stepik \enquote{Введение в Linux!}


\printbibliography



\end{document}
